\documentclass[12pt]{article}
\usepackage[utf8]{inputenc}
\usepackage[T1]{fontenc}
\usepackage{hyperref}
\title{Garza Language Resources for K-12 Curriculum Development}
\author{Compiled by Claude Sonnet 5.5 (claude-sonnet-5-5) for Andr\'as Kornai}
\date{}
\begin{document}
\maketitle

\section{Introduction}
Garza (Meak\'an; ISO 639-3 xgr) is an extinct Comecrudan language once spoken in Texas and Mexico. It is known only from two tribal names and twenty-one words recorded from the chief of the Garza by Jean-Louis Berlandier in 1828. By then the Garza spoke Spanish. Classification is disputed. Almost nothing exists, so a K-12 curriculum in the language is not feasible. At most, a heritage unit could be built around the history and the surviving record.

\section{Primary Educational Resources}
I found no primers, textbooks or school programmes.

\section{Linguistic and Reference Materials}
Wikipedia summarises the language: \url{https://en.wikipedia.org/wiki/Garza_language}. Its page on the Comecrudan family is at \url{https://en.wikipedia.org/wiki/Comecrudan_languages}. The page says Berlandier's manuscripts hold the only records of Garza and Mamulique.

\section{Religious and Community Resources}
I found no scripture translations or community organisations in the results.

\section{Digital and Audio Resources}
I found no audio, video or software resources.

\section{Curriculum Development Context}
The record is about twenty-one words, so nothing beyond a word list can be taught as the language itself. Any school use would be a history or linguistics unit. The sources seen are encyclopedic only. Berlandier's original manuscripts would need to be found through archives, and I did not see them in these results.
\end{document}